\documentclass[9pt,a4paper]{extarticle}

% ============================================================
% Formelsammlung – Allgemeine Elektrotechnik 1 (AET1)
% Layout: Kachel-Raster (tcolorbox), max. 2 A4-Seiten (Querformat)
%
% Hinweis: Das Raster wird HIER manuell mit fixer Breite pro Kachel
% gebaut (statt \tcbraster), weil die automatische tcolorbox-Raster-
% Umgebung im Test nicht zuverlässig mehrspaltig umgebrochen hat.
% ============================================================

\usepackage[a4paper,margin=5mm,landscape]{geometry}
\usepackage[utf8]{inputenc}
\usepackage[T1]{fontenc}
\usepackage[ngerman]{babel}
\usepackage{amsmath,amssymb}
\usepackage{siunitx}
\usepackage[skins]{tcolorbox}
\usepackage{enumitem}
\usepackage{xcolor}

\setlength{\parindent}{0pt}
\setlist{nosep, leftmargin=0.9em, itemsep=0pt, topsep=1pt}
\pagestyle{empty}

% Kurzbefehl für zentrierte Formel-Zeilen
\newcommand{\formel}[1]{\[#1\]}

% ---- Themenfarben --------------------------------------------------
\definecolor{cGrund}{HTML}{2C5AA0}   % Grundlagen / Gleichstrom
\definecolor{cNetz}{HTML}{1F8A70}    % Netzwerke / Quellen
\definecolor{cBauteil}{HTML}{B5651D} % Bauelemente C/L
\definecolor{cFeld}{HTML}{7B4EA3}    % Felder
\definecolor{cWechsel}{HTML}{C0392B} % Wechselstrom
\definecolor{cRef}{HTML}{555555}     % Konstanten / Einheiten

% ---- Manuelles 4-Spalten-Raster ------------------------------------
% Breite: (Textbreite 287mm - 3x Spaltenabstand 2mm) / 4 Spalten
\newlength{\tilew}
\setlength{\tilew}{69.25mm}
\newcommand{\tilegap}{\hspace{2mm}}

\tcbset{
  tilestyle/.style={
    nobeforeafter, valign=top, halign=left,
    width=\tilew, height=90mm,
    boxrule=0.4pt, arc=0pt,
    left=2pt, right=2pt, top=1.5pt, bottom=1.5pt, boxsep=0pt,
    fonttitle=\bfseries\scriptsize, coltitle=white,
    title
  }
}
% #1 Farbe  #2 Titel  #3 Inhalt
\newcommand{\tile}[3]{%
  \begin{tcolorbox}[tilestyle, colframe=#1, colback=#1!4, colbacktitle=#1, title={#2}]
    #3
  \end{tcolorbox}%
}

\begin{document}
\scriptsize

{\centering\bfseries\normalsize Formelsammlung AET1 -- Allgemeine Elektrotechnik 1\par}
\vspace{1mm}

% ============================================================
% SEITE 1 -- Gleichstromtechnik / Netzwerke
% TODO: Themen folgen noch -- Platzhalter mit typischen AET1-Inhalten,
% koennen ersetzt / ergaenzt / geloescht werden.
% ============================================================
\noindent
\tile{cGrund}{Grundgrößen}{%
  \begin{itemize}
    \item Ladung: $Q = I \cdot t$ \ [\si{\coulomb}]
    \item Strom: $I = \dfrac{Q}{t} = \dfrac{dQ}{dt}$
    \item Spannung: $U = \dfrac{W}{Q}$
    \item Leistung: $P = U \cdot I$
    \item Arbeit: $W = P \cdot t = U I t$
  \end{itemize}
}\tilegap
\tile{cGrund}{Ohmsches Gesetz}{%
  \formel{U = R I \quad R = \dfrac{U}{I} \quad I = \dfrac{U}{R}}
  \begin{itemize}
    \item Leitwert: $G = \dfrac{1}{R}$ \ [\si{\siemens}]
    \item $R = \rho \dfrac{l}{A} = \dfrac{l}{\kappa A}$
    \item $R(\vartheta) = R_0 (1+\alpha\Delta\vartheta)$
  \end{itemize}
}\tilegap
\tile{cGrund}{Kirchhoff}{%
  \textbf{Knotenregel:}
  \formel{\textstyle\sum I_{\text{ein}} = \sum I_{\text{aus}}}
  \textbf{Maschenregel:}
  \formel{\textstyle\sum_k U_k = 0}
}\tilegap
\tile{cNetz}{Widerstandsnetzwerke}{%
  Reihe: $R_{ges} = R_1+R_2+\dots$

  Parallel: $\dfrac{1}{R_{ges}} = \dfrac{1}{R_1}+\dfrac{1}{R_2}+\dots$

  2 Widerst.: $R_{ges}=\dfrac{R_1 R_2}{R_1+R_2}$
}

\vspace{3mm}\par\noindent
\tile{cNetz}{Spannungs-/Stromteiler}{%
  Spannungsteiler (unbelastet):
  \formel{U_2 = U \cdot \dfrac{R_2}{R_1+R_2}}
  Stromteiler:
  \formel{I_1 = I \cdot \dfrac{R_2}{R_1+R_2}}
}\tilegap
\tile{cNetz}{Leistung \& Energie}{%
  \begin{itemize}
    \item $P = UI = I^2R = \dfrac{U^2}{R}$
    \item Wirkungsgrad: $\eta = \dfrac{P_{ab}}{P_{zu}}$
    \item Leistungsanpassung: $R_L = R_i$
  \end{itemize}
}\tilegap
\tile{cNetz}{Quellen}{%
  Reale Spannungsquelle: \ $U = U_0 - R_i I$

  Reale Stromquelle: \ $I = I_0 - \dfrac{U}{R_i}$

  Umwandlung: $U_0 = I_0 R_i$
}\tilegap
\tile{cNetz}{Zweipoltheorie}{%
  \begin{itemize}
    \item \textbf{Thévenin:} $U_0=U_{Leerlauf}$, $R_i=R_{ab}$ (Quellen passiviert)
    \item \textbf{Norton:} $I_0=I_{Kurzschluss}$
  \end{itemize}
}

\newpage

% ============================================================
% SEITE 2 -- Bauelemente / Felder / Wechselstrom / Referenz
% ============================================================
\noindent
\tile{cBauteil}{Kondensator}{%
  \begin{itemize}
    \item $Q = CU$
    \item Platten-Kond.: $C = \varepsilon_0\varepsilon_r\dfrac{A}{d}$
    \item Energie: $W=\tfrac{1}{2}CU^2=\tfrac{1}{2}QU$
  \end{itemize}
  Reihe: $\dfrac{1}{C_{ges}}=\dfrac{1}{C_1}+\dfrac{1}{C_2}+\dots$

  Parallel: $C_{ges}=C_1+C_2+\dots$
}\tilegap
\tile{cBauteil}{RC-Vorgang}{%
  \formel{u_C(t) = U_0\!\left(1-e^{-t/\tau}\right)}
  \formel{\tau = RC}
}\tilegap
\tile{cBauteil}{Spule / Induktivität}{%
  \begin{itemize}
    \item $u_L = L\dfrac{di}{dt}$
    \item Energie: $W=\tfrac{1}{2}LI^2$
    \item Zylinderspule: $L=\mu_0\mu_r\dfrac{N^2A}{l}$
  \end{itemize}
  Reihe: $L_{ges}=L_1+L_2+\dots$

  Parallel: $\dfrac{1}{L_{ges}}=\dfrac{1}{L_1}+\dfrac{1}{L_2}+\dots$
}\tilegap
\tile{cBauteil}{RL-Vorgang}{%
  \formel{i_L(t) = I_0\!\left(1-e^{-t/\tau}\right)}
  \formel{\tau = \dfrac{L}{R}}
}

\vspace{3mm}\par\noindent
\tile{cFeld}{Elektrisches Feld}{%
  \begin{itemize}
    \item Feldstärke: $E = \dfrac{U}{d}$
    \item Kraft: $F = QE$
    \item Verschiebungsd.: $D=\varepsilon_0\varepsilon_r E$
  \end{itemize}
}\tilegap
\tile{cFeld}{Magnetisches Feld}{%
  \begin{itemize}
    \item Durchflutung: $\Theta = NI$
    \item Feldstärke: $H=\dfrac{NI}{l}$
    \item Flussdichte: $B=\mu_0\mu_r H$
    \item Fluss: $\Phi = BA$
    \item Induktion: $u_i=-N\dfrac{d\Phi}{dt}$
    \item Kraft: $F = BIl$
  \end{itemize}
}\tilegap
\tile{cWechsel}{Wechselstromgrößen}{%
  \begin{itemize}
    \item $u(t)=\hat u\sin(\omega t+\varphi)$
    \item $\omega = 2\pi f = \dfrac{2\pi}{T}$
    \item Effektivwert: $U_{eff}=\dfrac{\hat u}{\sqrt2}$
    \item Gleichrichtwert: $\bar u=\dfrac{2}{\pi}\hat u$
  \end{itemize}
}\tilegap
\tile{cRef}{Konstanten \& SI-Vorsätze}{%
  $\varepsilon_0=\SI{8.854e-12}{\farad\per\meter}$

  $\mu_0=4\pi\times10^{-7}\,\si{\henry\per\meter}$

  $e=\SI{1.602e-19}{\coulomb}$
  \vspace{2pt}

  \begin{tabular}{@{}llll@{}}
    p:$10^{-12}$ & n:$10^{-9}$ & µ:$10^{-6}$ & m:$10^{-3}$ \\
    k:$10^{3}$ & M:$10^{6}$ & G:$10^{9}$ &
  \end{tabular}
}

\end{document}
